% from ega0-1-fr.tex
\oldpage[0\textsubscript{I}]{66}

les conditions de (7.2.7)~; comme $A$ est la limite projective des $A_i$, on
voit que la prop. (7.2.7) donne la description de \emph{tous} les anneaux
adiques du type considéré (et en particulier de tous les anneaux
\emph{adiques noethériens}).
\end{remark}

\begin{example}[7.2.12]
\label{0.7.2.12-fr}
Soient $B$ un anneau, $\mathfrak J$ un idéal de $B$ tel que
$\mathfrak J/\mathfrak J^2$ soit un module de type fini sur
$B/\mathfrak J$ (ou sur $B$, ce qui revient au même)~; posons
$A=\varprojlim_n B/\mathfrak J^{n+1}$~; $A$ est le séparé complété de $B$
muni de la topologie $\mathfrak J$-préadique. Si
$A_n=B/\mathfrak J^{n+1}$, il est immédiat que les $A_n$ vérifient les
conditions de (7.2.7)~; donc $A$ est un anneau adique et si
$\overline{\mathfrak J}$ est l'adhérence dans $A$ de l'image canonique de
$\mathfrak J$, $\overline{\mathfrak J}$ est un idéal de définition de $A$,
$\overline{\mathfrak J^n}$ est l'adhérence de l'image canonique de
$\mathfrak J^n$, $A/\overline{\mathfrak J^n}$ s'identifie à
$B/\mathfrak J^n$ et
$\overline{\mathfrak J}/\overline{\mathfrak J^2}$ est isomorphe à
$\mathfrak J/\mathfrak J^2$ en tant que
$(A/\overline{\mathfrak J})$-module. De même, si $N$ est tel que
$N/\mathfrak JN$ soit un $B$-module de type fini, et si on pose
$M_i=N/\mathfrak J^{i+1}N$, $M=\varprojlim M_i$ est un $A$-module de type
fini, isomorphe au séparé complété de $N$ pour la topologie
$\mathfrak J$-préadique, $\overline{\mathfrak J^n}M$ s'identifie à
l'adhérence de l'image canonique de $\mathfrak J^nN$, et
$M/\overline{\mathfrak J^n}M$ à $N/\mathfrak J^nN$.
\end{example}

\subsection{Anneaux préadiques noethériens}
\label{subsection:0.7.3-fr}

\begin{env}[7.3.1]
\label{0.7.3.1-fr}
Soient $A$ un anneau, $\mathfrak J$ un idéal de $A$, $M$ un $A$-module~;
nous désignerons par
$\widehat A=\varprojlim_n A/\mathfrak J^n$ (resp.
$\widehat M=\varprojlim_u M/\mathfrak J^nM$) le séparé complété de $A$
(resp. $M$) pour la topologie $\mathfrak J$-préadique. Soit
\[
  M'\xrightarrow{u}M\xrightarrow{v}M''\to0
\]
une suite exacte de $A$-modules~; comme
$M/\mathfrak J^nM=M\otimes_A(A/\mathfrak J^n)$, la suite
\[
  M'/\mathfrak J^nM'\xrightarrow{u_n}
  M/\mathfrak J^nM\xrightarrow{v_n}
  M''/\mathfrak J^nM''\to0
\]
est exacte pour tout $n$. En outre, comme
$v(\mathfrak J^nM)=\mathfrak J^nv(M)=\mathfrak J^nM''$,
$\widehat v=\varprojlim v_n$ est surjective (Bourbaki,
\emph{Top. gén.}, chap.~IX, 2\textsuperscript{e} éd., p.~60, cor.~2).
D'autre part, si $z=(z_k)$ est un élément du noyau de $\widehat v$, pour tout
entier $k$, il existe un $z'_k\in M'/\mathfrak J^kM'$ tel que
$u_k(z'_k)=z_k$~; on en conclut qu'il existe
$z'=(z'_n)\in\widehat{M'}$ tel que les $k$ premières composantes de
$\widehat u(z')$ coïncident avec celles de $z$~; autrement dit, l'image par
$\widehat u$ de $\widehat{M'}$ est \emph{dense} dans le noyau de
$\widehat v$.

Si on suppose $A$ \emph{noethérien}, il en est de même de $\widehat A$, en
vertu de (7.2.12), $\mathfrak J/\mathfrak J^2$ étant alors un $A$-module de
type fini. En outre~:
\end{env}

\begin{namedtheorem}[7.3.2]{Théorème de Krull}
\label{0.7.3.2-fr}
Soient $A$ un anneau noethérien, $\mathfrak J$ un idéal de $A$, $M$ un
$A$-module de type fini, $M'$ un sous-module de $M$~; alors la topologie
induite sur $M'$ par la topologie $\mathfrak J$-préadique de $M$ est identique
à la topologie $\mathfrak J$-préadique de $M'$.
\end{namedtheorem}

Cela résulte aussitôt du

\begin{namedlemma}[7.3.2.1]{Lemme d'Artin-Rees}
\label{0.7.3.2.1-fr}
Sous les hypothèses de (7.3.2), il existe un entier $p$ tel que, pour
$n\geq p$, on ait
\[
  M'\cap\mathfrak J^nM=\mathfrak J^{n-p}
  (M'\cap\mathfrak J^pM).
\]
\end{namedlemma}

Pour la démonstration, voir ([I], p.~2-04).

